\documentclass[conference]{IEEEtran}
\usepackage[utf8]{inputenc}
\usepackage[T1]{fontenc}
\usepackage[spanish,es-tabla]{babel}
\usepackage{amsmath,amssymb,amsfonts}
\usepackage{booktabs}
\usepackage{array}
\usepackage{hyperref}
\usepackage{enumitem}
\usepackage{xcolor}
\usepackage{graphicx}
\usepackage{longtable}
\usepackage{multirow}
\usepackage{geometry}
\geometry{margin=1.7cm}

\title{Revisión técnica para un análisis multidimensional multimodal en la Plataforma Neurobiomecánica VCE}
\author{Documento de trabajo metodológico}
\date{}

\begin{document}
\maketitle

\begin{abstract}
Este documento revisa técnicas adicionales de procesamiento digital de señales, neuroimagen computacional, fusión multimodal y análisis multidimensional que pueden complementar una plataforma que integra EMG, dinamometría, TAC, resonancia estructural, difusión, tractografía de la vía corticoespinal, segmentación de 19 zonas y comparación longitudinal antes/después. La propuesta central es añadir una capa de integración por paciente que no reemplace los módulos actuales, sino que consolide sus salidas en un espacio común de características, regiones, modalidades y tiempo. A partir de esta representación se pueden aplicar modelos de factores latentes, descomposición conjunta/individual, fusión por grafos, tractometría, radiomía, análisis tensorial y métricas de discordancia entre modalidades.
\end{abstract}

\section{Planteamiento del problema}
La plataforma actual ya puede producir mediciones parciales por modalidad: actividad mioeléctrica, fuerza, morfometría, segmentaciones, mapas de correlación, tractografía y métricas anatómicas. El siguiente avance metodológico no consiste únicamente en extraer más variables aisladas, sino en construir un \textit{modelo multidimensional por paciente} que agrupe todos los estudios disponibles y permita responder preguntas integradas:

\begin{itemize}[leftmargin=*]
    \item ¿Qué regiones cambian de forma coherente entre neuroimagen, tractografía y mediciones periféricas?
    \item ¿Qué modalidad explica más variación en el cambio antes/después?
    \item ¿Existen zonas con deterioro anatómico pero compensación funcional?
    \item ¿La vía corticoespinal muestra cambios compatibles con la evolución de fuerza, EMG o masa muscular?
    \item ¿Qué variables parecen redundantes y cuáles aportan información complementaria?
\end{itemize}

El objetivo recomendado es crear una capa llamada, por ejemplo, \texttt{integracion\_multidimensional}, que reciba las salidas de los módulos ya existentes y genere: una matriz de características, un tensor paciente--tiempo--modalidad--región, un reporte de contribuciones por modalidad, mapas de cambio integrado y tablas Excel interpretables.

\section{Representación matemática base}
Para agrupar los estudios de un paciente, se propone organizar la información en un tensor multimodal:

\begin{equation}
\mathcal{X}_{p,t,m,r,k},
\end{equation}

donde $p$ representa el paciente, $t$ el momento de evaluación (Antes/Después), $m$ la modalidad (EMG, fuerza, TAC, T1, DWI, CST, zonas anatómicas), $r$ la región anatómica o canal funcional, y $k$ la característica extraída. Para análisis tabular, este tensor también puede desplegarse como una matriz:

\begin{equation}
\mathbf{X} \in \mathbb{R}^{n \times d},
\end{equation}

donde cada fila representa una unidad de análisis. Dependiendo del objetivo, una fila puede ser un paciente, un paciente-momento, una región, una región-momento o un segmento de tracto. Las columnas corresponden a características normalizadas.

Para comparación longitudinal se recomienda usar cambios relativos y diferencias normalizadas:

\begin{equation}
\Delta x = x_{Despues} - x_{Antes},
\end{equation}

\begin{equation}
\Delta_{\%}x = 100 \cdot \frac{x_{Despues} - x_{Antes}}{|x_{Antes}| + \epsilon},
\end{equation}

y asimetrías laterales:

\begin{equation}
A_{LR} = 100 \cdot \frac{x_{Derecha} - x_{Izquierda}}{\frac{1}{2}(x_{Derecha}+x_{Izquierda}) + \epsilon}.
\end{equation}

Estas formulaciones permiten integrar mediciones con unidades distintas, siempre que se normalicen por modalidad antes de fusionarlas.

\section{Banco de características por modalidad}
\subsection{EMG y señales neuromusculares}
Además de RMS, iEMG, frecuencia media/mediana y correlación, se recomiendan técnicas de análisis no estacionario y no lineal:

\begin{itemize}[leftmargin=*]
    \item \textbf{Transformada wavelet continua y coherencia wavelet:} permite estudiar cambios de energía tiempo--frecuencia y sincronización entre músculos o entre EMG y fuerza. Es útil cuando la activación muscular no es estacionaria o cuando los eventos tienen duración variable.
    \item \textbf{Wavelet packet y energía por bandas:} genera una descomposición más fina que la STFT y permite crear firmas espectrales por músculo, lado y etapa.
    \item \textbf{Hilbert--Huang Transform / EMD:} separa modos intrínsecos de oscilación y calcula frecuencia instantánea. Es útil en señales no lineales y no estacionarias.
    \item \textbf{Entropía muestral, entropía aproximada y complejidad multiescala:} cuantifican regularidad o variabilidad de activación neuromuscular.
    \item \textbf{Análisis de recurrencia:} permite identificar patrones repetitivos, estabilidad y transiciones en señales de movimiento o fuerza.
    \item \textbf{Acople EMG--fuerza:} puede modelarse mediante coherencia, correlación cruzada, fase relativa, retardos y modelos de regresión con desfase temporal.
\end{itemize}

Una característica importante sería construir un índice de acople electromecánico:

\begin{equation}
C_{EMG,F}(f,t) = \frac{|S_{EMG,F}(f,t)|^2}{S_{EMG}(f,t)S_F(f,t)},
\end{equation}

donde $S_{EMG,F}$ es el espectro cruzado local y $S_{EMG}$, $S_F$ son espectros de potencia. Esta métrica puede resumirse por bandas fisiológicamente relevantes.

\subsection{Dinamometría y cinética}
En fuerza se recomienda extraer no solo picos, sino dinámica de generación de fuerza:

\begin{itemize}[leftmargin=*]
    \item tasa de desarrollo de fuerza, pendiente máxima y tiempo al pico;
    \item impulso mecánico $I=\int F(t)dt$;
    \item simetría izquierda/derecha;
    \item estabilidad de la curva de fuerza mediante varianza local, entropía y área bajo la curva;
    \item acople con EMG usando retardos óptimos y coherencia tiempo--frecuencia.
\end{itemize}

\subsection{TAC, morfometría y radiomía muscular}
Para TAC y tejidos musculares/adiposos se recomienda añadir radiomía y textura. Además de volumen y área, pueden extraerse:

\begin{itemize}[leftmargin=*]
    \item distribución de unidades Hounsfield por músculo;
    \item percentiles de intensidad;
    \item textura GLCM: contraste, homogeneidad, energía y correlación;
    \item GLRLM/GLSZM: longitud de rachas, tamaño de zonas y heterogeneidad;
    \item forma 3D: elongación, compacidad, esfericidad, superficie y relación superficie/volumen;
    \item infiltración grasa intramuscular y subcutánea.
\end{itemize}

Estas variables ayudan a detectar cambios que no se observan solo con volumen. Dos regiones con igual volumen pueden diferir en calidad tisular, textura o infiltración grasa.

\subsection{T1, FreeSurfer y 19 zonas anatómicas}
Las 19 zonas pueden convertirse en una tabla longitudinal por región:

\begin{equation}
\mathbf{Z}_{r} = [V_r, I_r, \sigma_r, P_{10,r}, P_{90,r}, Corr_r, \Delta V_r, \Delta I_r],
\end{equation}

donde $V_r$ es volumen, $I_r$ intensidad media, $\sigma_r$ desviación de intensidad, $P_{10}$ y $P_{90}$ percentiles, $Corr_r$ la correlación local antes/después y $\Delta$ los cambios longitudinales.

Se recomienda agregar un índice de discordancia regional:

\begin{equation}
D_r = w_1 |z(\Delta V_r)| + w_2 |z(1-Corr_r)| + w_3 |z(\Delta I_r)|,
\end{equation}

que ordena zonas con mayor cambio multidimensional.

\subsection{DWI, tractografía y vía corticoespinal}
La vía corticoespinal no debería resumirse solamente con número total de streamlines. Se recomienda añadir tractometría o perfil a lo largo del tracto:

\begin{equation}
T(s) = [FA(s), MD(s), RD(s), AD(s), Density(s), Curvature(s)], \quad s \in [0,1],
\end{equation}

donde $s$ parametriza la trayectoria del tracto desde corteza motora hasta tronco encefálico. Esto permite ubicar si el cambio ocurre cerca de corteza, cápsula interna, mesencéfalo, puente o bulbo.

Salidas recomendadas:

\begin{itemize}[leftmargin=*]
    \item perfil de FA/MD/densidad en 100 nodos a lo largo de la CST;
    \item densidad de streamlines por segmento anatómico;
    \item curvatura, longitud media y dispersión espacial;
    \item proporción de fibras que atraviesan mesencéfalo, puente y bulbo;
    \item índice CST integrado.
\end{itemize}

\section{Modelos de integración multidimensional}
\subsection{PCA multibloque}
La PCA multibloque separa variables por modalidad y permite identificar ejes globales de variación sin que una modalidad domine solo por tener más columnas. Cada bloque se normaliza y se proyecta hacia componentes latentes compartidos. Puede usarse como primer modelo exploratorio.

\subsection{JIVE: variación conjunta e individual}
JIVE descompone la información en componentes compartidos entre bloques y componentes específicos de cada modalidad:

\begin{equation}
\mathbf{X}^{(m)} = \mathbf{J}^{(m)} + \mathbf{A}^{(m)} + \mathbf{E}^{(m)},
\end{equation}

donde $\mathbf{J}^{(m)}$ es la variación conjunta, $\mathbf{A}^{(m)}$ la variación individual de la modalidad y $\mathbf{E}^{(m)}$ el ruido residual. Para VCE, esto permitiría distinguir si un cambio antes/después es global del sistema motor o aparece solo en neuroimagen, solo en músculo o solo en fuerza.

\subsection{Factorización latente tipo MOFA}
Los modelos de factores latentes permiten integrar modalidades con diferente número de variables y datos faltantes:

\begin{equation}
\mathbf{X}^{(m)} \approx \mathbf{W}^{(m)}\mathbf{Z} + \mathbf{E}^{(m)},
\end{equation}

donde $\mathbf{Z}$ son factores latentes comunes del paciente y $\mathbf{W}^{(m)}$ son cargas por modalidad. En el proyecto, un factor podría representar recuperación motora periférica, otro integridad corticoespinal y otro cambio morfológico muscular.

\subsection{sPLS, CCA y DIABLO}
Los modelos PLS/CCA buscan variables de distintas modalidades que covarían. Su versión dispersa permite seleccionar variables relevantes. En VCE puede usarse para encontrar relaciones como:

\begin{itemize}[leftmargin=*]
    \item cambios de FA en CST asociados a cambios de fuerza;
    \item cambios de EMG asociados a volumen muscular;
    \item zonas cerebrales cuya correlación local se relaciona con simetría de fuerza.
\end{itemize}

\subsection{Similarity Network Fusion}
Cuando haya varios pacientes o varias unidades de análisis, se puede construir una red de similitud por modalidad y fusionarlas. Cada modalidad produce una matriz $S^{(m)}$ de similitud entre pacientes, regiones o evaluaciones. SNF genera una red integrada que resume la estructura común entre modalidades.

\subsection{Multiple Kernel Learning}
MKL permite representar cada modalidad mediante un kernel y aprender pesos de importancia:

\begin{equation}
K = \sum_{m=1}^{M} \beta_m K_m, \quad \beta_m \geq 0.
\end{equation}

En VCE, $K_m$ puede venir de EMG, fuerza, TAC, T1, DWI, CST o zonas. El peso $\beta_m$ sirve como medida de contribución relativa de cada modalidad para una tarea de clasificación o predicción.

\subsection{Análisis tensorial}
Si se organiza la información como paciente $\times$ tiempo $\times$ región $\times$ modalidad, se puede aplicar descomposición CP/Tucker:

\begin{equation}
\mathcal{X} \approx \sum_{q=1}^{Q} \lambda_q \mathbf{a}_q \circ \mathbf{b}_q \circ \mathbf{c}_q \circ \mathbf{d}_q,
\end{equation}

lo que permite detectar patrones compartidos de cambio: por ejemplo, un factor que activa ``Después'', regiones motoras, CST y fuerza.

\subsection{Grafos multimodales por paciente}
Un enfoque especialmente útil para la plataforma es crear un grafo por paciente. Los nodos serían regiones anatómicas, músculos, lados corporales, segmentos de CST y variables clínicas. Las aristas representarían correlación, coherencia, conectividad anatómica, acople EMG--fuerza o proximidad anatómica.

Métricas útiles:

\begin{itemize}[leftmargin=*]
    \item centralidad de regiones;
    \item comunidades funcionales/anatómicas;
    \item nodos con mayor cambio antes/después;
    \item aristas que aparecen o desaparecen;
    \item distancia de red entre corteza motora, CST y músculo.
\end{itemize}

\section{Propuesta concreta para la plataforma VCE}
\subsection{Nivel 1: extracción de características}
Crear una sección futura:

\begin{verbatim}
--sections features_multimodales
\end{verbatim}

Salidas:

\begin{itemize}[leftmargin=*]
    \item \texttt{features\_emg.xlsx}
    \item \texttt{features\_dinamometria.xlsx}
    \item \texttt{features\_tac.xlsx}
    \item \texttt{features\_zonas\_cerebrales.xlsx}
    \item \texttt{features\_cst\_tractometria.xlsx}
\end{itemize}

\subsection{Nivel 2: tabla maestra}
Crear una tabla maestra:

\begin{verbatim}
resultados/integracion_multidimensional/features_master.xlsx
\end{verbatim}

Columnas recomendadas:

\begin{itemize}[leftmargin=*]
    \item paciente;
    \item etapa;
    \item modalidad;
    \item región;
    \item lado;
    \item variable;
    \item valor;
    \item valor normalizado;
    \item cambio absoluto;
    \item cambio porcentual;
    \item QC/fuente.
\end{itemize}

\subsection{Nivel 3: modelos de integración}
Crear una sección:

\begin{verbatim}
--sections integracion_multidimensional
\end{verbatim}

Modelos iniciales recomendados:

\begin{enumerate}[leftmargin=*]
    \item PCA multibloque.
    \item JIVE o aproximación conjunta/individual.
    \item MOFA-like si hay datos faltantes.
    \item sPLS/CCA para pares de modalidades.
    \item grafo multimodal paciente-específico.
\end{enumerate}

\subsection{Nivel 4: reporte clínico/investigativo}
Salidas recomendadas:

\begin{itemize}[leftmargin=*]
    \item \texttt{reporte\_integrado\_paciente.xlsx}
    \item \texttt{ranking\_zonas\_cambio.xlsx}
    \item \texttt{contribucion\_modalidades.xlsx}
    \item \texttt{grafo\_multimodal.graphml}
    \item \texttt{mapa\_cambio\_integrado.nii.gz}
    \item \texttt{dashboard\_paciente.html}
\end{itemize}

\section{Qué técnicas son más viables según el tamaño actual de datos}
Con pocos pacientes y pocas evaluaciones, no conviene empezar por modelos profundos ni clasificadores complejos. El orden recomendado es:

\begin{enumerate}[leftmargin=*]
    \item \textbf{Muy viable ahora:} tabla maestra, normalización, índices de cambio, tractometría, radiomía, PCA multibloque, correlaciones cruzadas y grafos descriptivos.
    \item \textbf{Viable con cuidado:} CCA/sPLS, JIVE, MOFA-like, SNF sobre regiones o pacientes-momento, análisis tensorial simple.
    \item \textbf{Para más adelante:} MKL supervisado, GNN, deep learning multimodal, modelos predictivos clínicos robustos.
\end{enumerate}

\section{Índices integrados propuestos}
\subsection{Índice de recuperación neuromotora}

\begin{equation}
IRN = \alpha_1 z(\Delta Fuerza) + \alpha_2 z(\Delta EMG) + \alpha_3 z(\Delta CST) + \alpha_4 z(\Delta Musculo) - \alpha_5 z(Deterioro),
\end{equation}

con pesos ajustables según criterio metodológico.

\subsection{Índice de discordancia central-periférica}

\begin{equation}
IDCP = |z(\Delta Neuroimagen) - z(\Delta Periferico)|.
\end{equation}

Este índice identifica pacientes o regiones donde el sistema central y periférico no evolucionan de forma concordante.

\subsection{Índice de integridad corticoespinal segmentaria}

\begin{equation}
I_{CST}(s)= w_1 z(FA(s)) - w_2 z(MD(s)) + w_3 z(Density(s)) - w_4 z(Curvature(s)).
\end{equation}

Permite comparar segmentos específicos de la CST entre Antes y Después.

\section{Conclusión técnica}
Sí es posible hacer un análisis multidimensional que agrupe todos los estudios de un paciente. La recomendación principal es no empezar por un modelo único demasiado complejo, sino construir una arquitectura escalonada: extracción de características por modalidad, tabla maestra longitudinal, normalización, modelos exploratorios multibloque, grafos por paciente y reportes de contribución por modalidad. Las técnicas con mejor relación entre potencia interpretativa y viabilidad actual son tractometría CST, radiomía por ROI, PCA multibloque, JIVE, MOFA-like, sPLS/CCA, SNF y grafos multimodales. Esta capa permitiría pasar de reportes separados por estudio a un perfil integrado de recuperación neuromotora.

\bibliographystyle{IEEEtran}
\bibliography{referencias}

\end{document}
